\documentclass[12pt]{article}
\usepackage[T1]{fontenc}
\usepackage[utf8]{inputenc}
\usepackage{lmodern}
\usepackage{textcomp}
\usepackage{enumitem}
\usepackage{geometry}
\usepackage{graphicx}
\usepackage{amsmath, amssymb}
\geometry{margin=1in}
\begin{document}
\begin{center}
Physics - Undergraduate Level Assessment 18 \
Total Marks: 40 \
Answer all questions.
\end{center}

\section*{Multiple Choice (10 marks)}
\begin{enumerate}[label=\arabic*.]
\item The eigenvalues of a Hermitian operator are always
\begin{enumerate}[label=(\alph*)]
    \item real
    \item imaginary
    \item degenerate
    \item linear
\end{enumerate}
\item The emission spectrum of the doubly ionized lithium atom Li++ (Z = 3, A = 7) is identical to that of a hydrogen atom in which all the wavelengths are
\begin{enumerate}[label=(\alph*)]
    \item decreased by a factor of 9
    \item decreased by a factor of 49
    \item decreased by a factor of 81
    \item increased by a factor of 9
\end{enumerate}
\item A three-dimensional harmonic oscillator is in thermal equilibrium with a temperature reservoir at temperature T. The average total energy of the oscillator is
\begin{enumerate}[label=(\alph*)]
    \item (1/2) k T
    \item kT
    \item (3/2) k T
    \item 3 kT
\end{enumerate}
\item For which of the following thermodynamic processes is the increase in the internal energy of an ideal gas equal to the heat added to the gas?
\begin{enumerate}[label=(\alph*)]
    \item Constant temperature
    \item Constant volume
    \item Constant pressure
    \item Adiabatic
\end{enumerate}
\item An organ pipe, closed at one end and open at the other, is designed to have a fundamental frequency of C (131 Hz). What is the frequency of the next higher harmonic for this pipe?
\begin{enumerate}[label=(\alph*)]
    \item 44 Hz
    \item 196 Hz
    \item 262 Hz
    \item 393 Hz
\end{enumerate}
\end{enumerate}

\section*{True / False (10 marks)}
\textit{Answer True or False}
\begin{enumerate}[label=\arabic*.]
\setcounter{enumi}{5}
\item True or False: In a single-electron atom with l = 2, there are only 4 allowed values for the magnetic quantum number m\_l.
\item True or False: Electromagnetic radiation emitted from a nucleus is most likely to be in the form of gamma rays.
\item True or False: A dye laser is the most suitable type of laser for spectroscopy across a range of visible wavelengths.
\item True or False: A particle that decays in 2.0 ms in its rest frame will travel 360 m in the lab frame at a speed of 0.60 c.
\item True or False: In a reversible thermodynamic process, the entropy of the system and its environment remains unchanged.
\end{enumerate}

\section*{Long Form Response (20 marks)}
\textit{Answer each question with a clear and complete explanation.}
\begin{enumerate}[label=\arabic*.]
\setcounter{enumi}{10}
\item Discuss the conditions under which a beam of electrons can diffract on a crystal surface, and analyze the relationship between the kinetic energy of the electrons and the observed scattering pattern, considering a lattice spacing of 0.4 nm. What is the approximate kinetic energy required to achieve this diffraction?
\item Discuss the relationship hypothesized by De Broglie between the linear momentum and wavelength of a free massive particle, emphasizing the role of Planck's constant in this context.
\end{enumerate}

\vfill
\noindent\textit{End of Paper}
\end{document}